%% ===========================================================================
%%  Project brief: paired systematic review + empirical companion study
%%  Prepared for Dr. Aminat Abiola Ajibola
%%  Doubles as (a) a research protocol and (b) a kickoff spec for Claude Code
%% ===========================================================================
\documentclass[11pt,a4paper]{article}

\usepackage[T1]{fontenc}
\usepackage[utf8]{inputenc}
\usepackage{charter}
\usepackage[scaled=0.88]{helvet}
\usepackage{microtype}
\usepackage[a4paper,margin=2.2cm,top=2.4cm,bottom=2.4cm]{geometry}
\usepackage{xcolor}
\usepackage{enumitem}
\usepackage{titlesec}
\usepackage{fancyhdr}
\usepackage{tcolorbox}
\usepackage{booktabs}
\usepackage{longtable}
\usepackage{array}
\usepackage{listings}
\usepackage[hidelinks]{hyperref}

\tcbuselibrary{skins,breakable}

\definecolor{ink}{HTML}{1A1A1A}
\definecolor{accent}{HTML}{8C2F1E}
\definecolor{rule}{HTML}{C9C2B8}
\definecolor{panelbg}{HTML}{F4F1EC}
\definecolor{softgrey}{HTML}{5C5851}
\definecolor{codebg}{HTML}{F7F5F2}
\definecolor{codekey}{HTML}{1F4E5F}

\color{ink}

\titleformat{\section}
  {\normalfont\sffamily\bfseries\large\color{accent}}{\thesection.}{0.6em}{}
\titleformat{\subsection}
  {\normalfont\sffamily\bfseries\normalsize\color{ink}}{\thesubsection}{0.6em}{}
\titleformat{\subsubsection}
  {\normalfont\sffamily\bfseries\small\color{softgrey}}{\thesubsubsection}{0.6em}{}
\titlespacing*{\section}{0pt}{2.4ex plus 1ex}{1.1ex}
\titlespacing*{\subsection}{0pt}{1.9ex plus 1ex}{0.7ex}

\setlist[enumerate,1]{leftmargin=1.7em,itemsep=0.45em,topsep=0.5em}
\setlist[itemize,1]{leftmargin=1.4em,itemsep=0.3em,topsep=0.4em,
  label=\textcolor{accent}{\textbullet}}

\pagestyle{fancy}
\fancyhf{}
\renewcommand{\headrulewidth}{0.4pt}
\fancyhead[L]{\footnotesize\sffamily\color{softgrey}Paired review + empirical study}
\fancyhead[R]{\footnotesize\sffamily\color{softgrey}Project brief}
\fancyfoot[C]{\footnotesize\sffamily\color{softgrey}\thepage}

\setlength{\parskip}{0.55em}
\setlength{\parindent}{0pt}

\newtcolorbox{panel}{breakable,enhanced,colback=panelbg,colframe=rule,
  boxrule=0.4pt,left=10pt,right=10pt,top=8pt,bottom=8pt,arc=1pt,
  before skip=10pt,after skip=12pt}
\newtcolorbox{keybox}{breakable,enhanced,colback=white,colframe=accent,
  boxrule=0.9pt,left=10pt,right=10pt,top=8pt,bottom=8pt,arc=1pt,
  before skip=10pt,after skip=12pt}

\lstdefinestyle{plain}{
  basicstyle=\ttfamily\scriptsize,
  backgroundcolor=\color{codebg},
  frame=single, rulecolor=\color{rule},
  breaklines=true, breakatwhitespace=false,
  columns=fullflexible, keepspaces=true,
  showstringspaces=false,
  xleftmargin=6pt, xrightmargin=0pt,
  framexleftmargin=6pt,
  aboveskip=10pt, belowskip=10pt,
  literate={─}{{-}}1 {│}{{|}}1 {├}{{|}}1 {└}{{\textbackslash}}1
}
\lstset{style=plain}

\begin{document}

\begin{center}
  {\sffamily\footnotesize\color{softgrey}\MakeUppercase{Project brief}}\\[0.9em]
  {\sffamily\LARGE\bfseries Machine Learning for Mental Health Prediction:\\[0.3em]
   A Systematic Appraisal and an Empirical Demonstration}\\[0.7em]
  {\sffamily\normalsize\color{softgrey} Paired-paper design, dataset plan, and
   repository specification}\\[0.9em]
  {\footnotesize\color{softgrey} Prepared for Dr.\ Aminat Abiola Ajibola \quad\textbullet\quad
   Includes a kickoff prompt for Claude Code}
\end{center}
\vspace{0.3em}
{\color{rule}\hrule height 0.6pt}
\vspace{1.2em}

% =====================================================================
\section{The argument in one page}

\begin{keybox}
\textbf{Paper A (review)} establishes that the field systematically fails to
externally validate, fails to report calibration, and frequently reports
performance inflated by subject-level data leakage.

\smallskip
\textbf{Paper B (empirical)} takes public mental-health datasets and
\emph{quantifies} each of those three failures on real data --- showing how
much accuracy inflates under leaky splitting, how far performance falls
under genuine external validation, and how badly miscalibrated the resulting
models are.

\smallskip
Neither paper is remarkable alone. The review restates a concern that
methodologists already voice; the empirical study is one more model on
familiar datasets. \textbf{Together they are a complete argument}: the review
measures how widespread the problem is, and the companion measures how much
it costs. That pairing is uncommon and is what makes this Q1-competitive.
\end{keybox}

\subsection{Why not the review alone}

Four directly competing reviews already exist, the most recent published in
January 2026 (Hasan et al., \emph{Frontiers in Digital Health}, 3{,}320
records screened). A descriptive review at this scope will be rejected on
novelty. Reframing to a methodological appraisal fixes that --- but an
appraisal that only counts deficiencies invites the reply
\emph{``so what?''} The companion study answers it with numbers.

\subsection{Why not the empirical study alone}

Primary ML-on-mental-health-data papers are abundant and the marginal one is
not publishable at Q1. What makes Paper~B defensible is that its research
questions are \emph{derived from} Paper~A rather than chosen arbitrarily.
The review licenses the experiment.

% =====================================================================
\section{Paper A --- the systematic review}

\subsection{Research question}

How methodologically sound and completely reported are machine learning
models developed to predict mental health disorders across populations?
A systematic review and critical appraisal using PROBAST and TRIPOD+AI.

\subsection{Design summary}

\begin{itemize}
  \item \textbf{Registration:} PROSPERO, before the searches are run.
  \item \textbf{Databases:} PubMed, Scopus, Web of Science, IEEE Xplore.
        PsycINFO if access can be obtained; declare its absence if not.
  \item \textbf{Window:} 2018--2026, English, peer-reviewed.
  \item \textbf{Population:} any, any setting. \textbf{Condition:} any mental
        health disorder. \textbf{Index model:} supervised ML/DL predicting a
        mental health outcome.
  \item \textbf{Appraisal:} PROBAST (20 signalling questions, 4 domains) and
        TRIPOD+AI adherence, dual independent reviewers, Cohen's $\kappa$.
  \item \textbf{Expected yield:} 5{,}000--9{,}000 records after
        deduplication; 80--200 included studies.
\end{itemize}

\subsection{Primary outcomes --- these become Paper B's hypotheses}

\begin{enumerate}
  \item Proportion of studies performing \textbf{any external validation}
  \item Proportion reporting \textbf{calibration}
  \item Proportion reporting \textbf{AUC} rather than accuracy alone
  \item Proportion with \textbf{leakage-prone evaluation design} --- notably
        record-wise rather than subject-wise splitting of repeated measures
  \item Proportion meeting the \textbf{EPV $\geq$ 200} expectation for ML
        models (Moons et al., 2025)
  \item Proportion releasing \textbf{code} or \textbf{data}
\end{enumerate}

Full protocol, PICOS criteria table and the four database search strings are
in \texttt{P1\_review\_protocol.md}. The extraction and appraisal workbook is
\texttt{P1\_extraction\_appraisal.xlsx}.

% =====================================================================
\section{Paper B --- the empirical companion}

\subsection{Three experiments, each answering one review finding}

\begin{panel}
\textbf{E1 --- The cost of leaky splitting.}
Train the same model on the same data under (a) record-wise splitting, where
windows from one participant may appear in both train and test, and
(b) subject-wise splitting, where no participant crosses the boundary.
Report the difference.

\smallskip
\emph{Why this matters:} actigraphy and sensor studies produce many windows
per participant. Splitting by window rather than by person is a common and
often unremarked practice, and it inflates accuracy substantially. E1 puts a
number on that inflation.

\medskip
\textbf{E2 --- The cost of skipping external validation.}
Train on cohort~1, evaluate unchanged on cohort~2 with no refitting.
Contrast with the internally cross-validated estimate on cohort~1.

\smallskip
\emph{Why this matters:} the review will show most studies report only
internal validation. E2 measures the optimism that entails.

\medskip
\textbf{E3 --- The calibration blind spot.}
For every model in E1 and E2, report calibration slope, intercept, Brier
score and a reliability curve alongside discrimination.

\smallskip
\emph{Why this matters:} a model can discriminate well and still produce
probability estimates that are useless for clinical decisions. The review
will show this is almost never reported.
\end{panel}

\subsection{Datasets}

\begin{center}
\footnotesize
\setlength{\tabcolsep}{5pt}
\renewcommand{\arraystretch}{1.3}
\begin{tabular}{@{}p{0.15\linewidth}p{0.33\linewidth}p{0.14\linewidth}p{0.29\linewidth}@{}}
\toprule
\textbf{Dataset} & \textbf{Content} & \textbf{Access} & \textbf{Link}\\
\midrule
Depresjon & Actigraphy; 23 unipolar/bipolar depressed + 32 controls; MADRS &
  Open, cite req. & \url{https://datasets.simula.no/depresjon/}\\
Psykose & Actigraphy; 22 schizophrenia + 32 controls; BPRS &
  Open, cite req. & \url{https://datasets.simula.no/psykose/}\\
HYPERAKTIV & Actigraphy; ADHD patients &
  Open, cite req. & \url{https://datasets.simula.no/hyperaktiv/}\\
OBF-Psychiatric & Aggregates the three above; \emph{Sci Data} 2025 &
  Open & \url{https://www.nature.com/articles/s41597-025-04384-3}\\
DAIC-WOZ / E-DAIC & 189 clinical interviews; audio, transcripts, facial
  features; PHQ-8 & EULA required &
  \url{https://dcapswoz.ict.usc.edu/}\\
\bottomrule
\end{tabular}
\end{center}

\smallskip
Simula dataset index: \url{https://github.com/simula/datasets.simula.no}

\subsubsection*{Candidates to verify before committing}

Named here for completeness; confirm access terms before building against
them. \textbf{NHANES} (CDC; PHQ-9 depression module; large, population
representative, fully open). \textbf{StudentLife} (Dartmouth; smartphone
sensing in university students). \textbf{MODMA} (Lanzhou University; EEG,
depression). \textbf{WESAD} (wearable stress). \textbf{CLPsych} and
\textbf{eRisk} shared-task corpora (social media text; signed agreements
required, and note the ethical debate around scraped clinical inference).

\subsection{The design rule that makes this work}

\begin{keybox}
\textbf{Do not pool the datasets into one training corpus.}

\smallskip
Pooling Depresjon, Psykose, DAIC and NHANES and reporting cross-validated
accuracy would commit exactly the error Paper~A criticises: different
populations, instruments, label definitions and prevalences, silently merged.
A reviewer who has read your own review would reject it.

\smallskip
\textbf{Use them as separate cohorts.} Train on one, validate unchanged on
another. The actigraphy family is well suited because the modality is
constant while the populations differ, so the comparison isolates cohort
shift rather than measurement shift.

\smallskip
The expected result --- substantial performance degradation out of sample,
with calibration degrading faster than discrimination --- \emph{is} the
paper.
\end{keybox}

\subsection{Reporting standard}

Paper~B follows \textbf{TRIPOD+AI} as an author, not as an appraiser. Every
item the review scores other studies against must be satisfied here. Anything
less would be indefensible given the paper's argument.

% =====================================================================
\section{Repository specification}

\subsection{Proposed structure}

\begin{lstlisting}
mlmh-appraisal/
├── README.md
├── pyproject.toml              # uv or poetry; pin every version
├── LICENSE                     # MIT for code
├── .gitignore                  # MUST exclude data/raw/ entirely
├── configs/
│   ├── base.yaml               # seeds, CV scheme, metrics
│   ├── e1_split_leakage.yaml
│   ├── e2_external_val.yaml
│   └── e3_calibration.yaml
├── data/
│   ├── raw/                    # gitignored. Never committed.
│   ├── interim/                # gitignored
│   ├── processed/              # gitignored
│   └── README.md               # how to obtain each dataset + checksums
├── src/mlmh/
│   ├── data/
│   │   ├── loaders.py          # one loader per dataset -> common schema
│   │   ├── schema.py           # subject_id, window_id, features, label, cohort
│   │   └── windowing.py        # segmentation; carries subject_id through
│   ├── features/
│   │   └── actigraphy.py       # statistical + circadian features
│   ├── models/
│   │   ├── registry.py         # baseline, logreg, RF, XGBoost, 1D-CNN
│   │   └── pipelines.py        # sklearn Pipelines only; no manual prefit
│   ├── evaluation/
│   │   ├── splitters.py        # SubjectWiseSplit vs RecordWiseSplit
│   │   ├── metrics.py          # AUC, macro-F1, Brier, calib slope/intercept
│   │   ├── bootstrap.py        # BCa CIs, subject-level resampling
│   │   └── external.py         # train-on-A, test-on-B runner
│   ├── reporting/
│   │   ├── tables.py           # emits LaTeX tables directly
│   │   ├── figures.py          # reliability curves, forest plots
│   │   └── tripod.py           # TRIPOD+AI self-audit checklist
│   └── cli.py                  # entry point: python -m mlmh run <config>
├── scripts/
│   ├── prepare_data.py
│   ├── run_e1.py
│   ├── run_e2.py
│   └── run_e3.py
├── tests/
│   ├── test_no_subject_leakage.py   # THE critical test; see below
│   ├── test_pipeline_fit_order.py
│   └── test_metrics.py
├── results/                    # committed: tables, figures, run manifests
└── paper/
    ├── review/                 # Paper A sources
    └── empirical/              # Paper B sources
\end{lstlisting}

\subsection{Non-negotiable engineering rules}

These exist because the paper's credibility depends on them. Each should be
enforced by a test, not by discipline.

\begin{enumerate}

  \item \textbf{Subject-wise splitting is the default everywhere.} Record-wise
        splitting exists \emph{only} inside the E1 experiment, where it is the
        object of study. \texttt{test\_no\_subject\_leakage.py} must assert
        that no \texttt{subject\_id} appears in both train and test index sets
        for any fold of any non-E1 run, and must fail the build if one does.

  \item \textbf{All preprocessing lives inside the \texttt{Pipeline}.}
        Scaling, imputation and any resampling are fitted on training folds
        only. No \texttt{fit\_transform} on the full dataset before splitting,
        ever. \texttt{test\_pipeline\_fit\_order.py} guards this.

  \item \textbf{Resampling (SMOTE or otherwise) inside CV folds only},
        using \texttt{imblearn.pipeline.Pipeline} rather than the sklearn one,
        which will silently resample the validation fold.

  \item \textbf{Bootstrap at the subject level, not the window level.}
        Resampling windows treats correlated observations as independent and
        produces confidence intervals that are far too narrow.

  \item \textbf{Every run writes a manifest} --- git SHA, config hash, seed,
        library versions, dataset checksums --- into \texttt{results/}.
        A results table that cannot be traced to a manifest does not go in the
        paper.

  \item \textbf{Seeds fixed and recorded.} Report across at least 10 seeds,
        not one.

  \item \textbf{Raw data never enters git.} Redistribution terms differ per
        dataset, and DAIC-WOZ's EULA prohibits it outright.
        \texttt{data/README.md} carries download instructions and SHA-256
        checksums instead.

\end{enumerate}

\subsection{Reporting outputs the code should generate directly}

Have the pipeline emit publication-ready artefacts rather than numbers to be
transcribed by hand --- transcription is where errors enter.

\begin{itemize}
  \item LaTeX tables written straight into \texttt{paper/empirical/tables/}
  \item Reliability curves and discrimination/calibration comparison figures
        as PDF
  \item A TRIPOD+AI self-audit report listing which items the manuscript
        satisfies, generated from the run manifests
\end{itemize}

% =====================================================================
\section{Kickoff prompt for Claude Code}

Paste the block below into a new Claude Code session in an empty repository.
It states the scientific constraints first, because those are what the
implementation must serve.

\begin{lstlisting}
I am building the codebase for an empirical companion study to a systematic
review of machine learning models for mental health prediction. The review
found that published models rarely perform external validation, rarely report
calibration, and frequently use evaluation designs that leak subject
information. This repository must demonstrate and quantify those three
failures rigorously -- so its own methodology has to be beyond criticism.

Please scaffold a Python project with this structure: [paste the tree from
Section 4.1 of the brief]

Core scientific constraints, which override convenience:

1. SUBJECT-WISE SPLITTING IS THE DEFAULT. Actigraphy data yields many windows
   per participant. No subject_id may appear in both train and test in any
   fold. Implement SubjectWiseSplit and RecordWiseSplit as separate splitter
   classes. RecordWiseSplit is used ONLY in experiment E1, where the
   inflation it causes is what we are measuring.

2. Write tests/test_no_subject_leakage.py FIRST, before any model code. It
   must fail the build if any non-E1 fold shares a subject between train and
   test.

3. All preprocessing goes inside an imblearn Pipeline so that scaling,
   imputation and resampling are fitted on training folds only. Never
   fit_transform the full dataset before splitting.

4. Bootstrap confidence intervals must resample at the SUBJECT level, not
   the window level.

5. Every run writes a manifest to results/ containing git SHA, config hash,
   seed, library versions and dataset checksums.

6. data/raw/ is gitignored. Never commit dataset files -- DAIC-WOZ's EULA
   prohibits redistribution.

Three experiments to support:
  E1: record-wise vs subject-wise splitting on the same data; report the
      accuracy/AUC inflation.
  E2: train on cohort A, evaluate unchanged on cohort B (no refitting);
      contrast with the internal CV estimate on A.
  E3: for every model above, report calibration slope, intercept, Brier
      score and reliability curves alongside discrimination.

Datasets (I will download and place these in data/raw/ myself):
  - Depresjon:  https://datasets.simula.no/depresjon/
  - Psykose:    https://datasets.simula.no/psykose/
  - HYPERAKTIV: https://datasets.simula.no/hyperaktiv/
  - DAIC-WOZ:   https://dcapswoz.ict.usc.edu/  (EULA pending)
Each is per-subject CSV of minute-level activity counts plus a scores file
with demographics and clinical scales. Write one loader per dataset mapping
into a common schema: subject_id, cohort, window_id, features, label.

Models: majority-class baseline, logistic regression, random forest,
XGBoost, and a small 1D-CNN. The first two exist to establish that the
others earn their complexity.

Metrics: macro-F1, AUROC (one-vs-rest where multiclass), Brier score,
calibration slope and intercept, with BCa bootstrap CIs.

Start by scaffolding the repo, writing pyproject.toml with pinned versions,
and implementing schema.py, splitters.py and the leakage test. Do not write
model code until the leakage test passes against synthetic fixture data.
\end{lstlisting}

\subsection{Why this order}

Asking Claude Code to build the leakage test before the models is
deliberate. If the splitter is wrong, every downstream number is wrong and
the error is invisible --- results simply look better than they should. A
failing test at the outset is much cheaper than a retraction.

% =====================================================================
\section{Sequencing across both papers}

\begin{enumerate}
  \item \textbf{Now.} Register Paper~A on PROSPERO. Submit the DAIC-WOZ
        EULA request --- it takes time and is the only external dependency.
        Download the three Simula datasets, which are immediately available.
  \item \textbf{Weeks 1--2.} Scaffold the repository. Get the leakage test
        passing against synthetic fixtures. Run the four database searches
        for Paper~A and record the counts.
  \item \textbf{Weeks 3--8.} Screening and extraction for Paper~A, in Rayyan
        and the extraction workbook. In parallel, implement loaders,
        features and the model registry for Paper~B.
  \item \textbf{Weeks 8--12.} Run E1--E3. Paper~A's appraisal results
        determine which findings Paper~B should foreground.
  \item \textbf{Weeks 12--16.} Draft both manuscripts. Paper~A cites Paper~B
        as a demonstration; Paper~B cites Paper~A as motivation.
\end{enumerate}

\subsection{Submission strategy}

Submit \textbf{Paper~A first}, or both simultaneously to different journals.
Do not submit Paper~B alone --- without the review establishing that the
problem is widespread, its contribution reads as a methodological remark
rather than a finding.

Cross-cite explicitly. If Paper~A is still under review when Paper~B is
submitted, cite it as \emph{under review} with a preprint DOI, and say so
in the cover letter. Editors accept this readily; discovering it later
independently does not go as well.

\subsection{Candidate venues}

\begin{center}
\footnotesize
\setlength{\tabcolsep}{6pt}
\renewcommand{\arraystretch}{1.25}
\begin{tabular}{@{}p{0.12\linewidth}p{0.40\linewidth}p{0.42\linewidth}@{}}
\toprule
& \textbf{Paper A (review)} & \textbf{Paper B (empirical)}\\
\midrule
Primary & Journal of Biomedical Informatics & Journal of Biomedical Informatics\\
Alt.\ 1 & Int.\ Journal of Medical Informatics & Computers in Biology and Medicine\\
Alt.\ 2 & JMIR Mental Health & Artificial Intelligence in Medicine\\
\bottomrule
\end{tabular}
\end{center}

Verify current quartile rankings in JCR or Scopus before committing. Avoid
sending both to the same journal simultaneously without telling the editor.

% =====================================================================
\section{Governance}

\begin{itemize}
  \item \textbf{DAIC-WOZ EULA} --- start now; it gates any use of that cohort.
  \item \textbf{Simula datasets} --- open for research and education, but the
        source publications must be cited. Commercial use needs written
        permission.
  \item \textbf{IRB} --- even for public de-identified data, check whether
        Hafr Al-Batin requires a secondary-use determination. Usually a short
        expedited form.
  \item \textbf{Code licence} --- MIT for the repository. Do not apply a code
        licence to data you did not create.
  \item \textbf{Social media corpora} --- if CLPsych or eRisk are added later,
        note that inferring clinical status from scraped posts carries ethical
        considerations that reviewers at clinical venues will raise. Address
        it in the manuscript rather than waiting to be asked.
\end{itemize}

\vspace{1em}
{\footnotesize\color{softgrey}
Companion files: \texttt{P1\_review\_protocol.md} (PICOS criteria, four
database search strings, PROSPERO content) and
\texttt{P1\_extraction\_appraisal.xlsx} (PROBAST and TRIPOD+AI extraction
workbook, 7 sheets).\par}

\end{document}
